\section{Structural rationality: exact formal boundaries}
\label{sec:p11-rat105-annex}

These proofs concern the matrix at printed 626 and the elementary
expected-value compatibility discussed at printed 629--630
\parencite{NidaRumelinGutwaldZuber2021}. They establish no philosophical
identity between rationality, morality, efficiency, and conformity.

\begin{annexitem}{annex:p11.rat105.pd}{thm:p11-rat105-pd}{Dominance and every equilibrium of the source matrix}
\begin{proof}
Against cooperation, either player receives 3 from cooperating and 4
from defecting. Against defection, the values are 1 and 2. If the
opponent's cooperation probability is \(q\in[0,1]\), expected utilities
are
\[
 3q+1(1-q)=1+2q,\qquad
 4q+2(1-q)=2+2q.
\]
Their difference is one at every \(q\), including both endpoints.
Thus the player's unique best response to every opponent mixture is
pure defection. An own mixture putting positive mass on cooperation
has strictly smaller expected utility than pure defection by that mass,
so it cannot be an equilibrium response. Consequently every independent
mixed equilibrium has both players defecting. The profile \((D,D)\)
is an equilibrium because both players are best responding there.
This proves existence and uniqueness for this matrix without invoking
a general equilibrium theorem or assumptions about higher-order beliefs.

Mutual defection gives \((2,2)\), whereas mutual cooperation gives
\((3,3)\); the latter improves each coordinate by one. This comparison
of joint outcomes does not alter the unilateral-deviation comparison.
In particular, at \(q=1\), the values of cooperation and defection
are 3 and 4: believing the other will cooperate is not enough to reverse
the supplied ranking.
\end{proof}
\end{annexitem}

\begin{annexitem}{annex:p11.rat105.value}{thm:p11-rat105-value}{Expected-value comparisons, continuity, scale and a structural example}
\begin{proof}
Every probability vector on finite nonempty \(Z\) gives a finite real
sum \(V_v(p)\). Completeness and transitivity of the induced weak
preference follow from those of the real order, including ties. For
\(0<\alpha<1\) and lotteries \(p,r,s\), finite distributivity gives
\[
 V_v(\alpha p+(1-\alpha)s)
   -V_v(\alpha r+(1-\alpha)s)
   =\alpha\bigl(V_v(p)-V_v(r)\bigr).
\]
Since \(\alpha>0\), the difference is nonnegative exactly when
the original difference is nonnegative. This is precisely mixture
independence, not probabilistic independence of events.

For the stated continuity condition, put
\(M=\sum_{z\in Z}(1+v(z)^2)>0\). Each coordinate difference
\(p'(z)-p(z)\) lies between \(-d(p',p)\) and \(d(p',p)\),
since its square is one term of the squared Euclidean distance.
Also \(-(1+v(z)^2)\leq v(z)\leq1+v(z)^2\), as follows by
completing the squares in \(v(z)^2+v(z)+1\) and
\(v(z)^2-v(z)+1\). Multiplying the two bounds gives
\[
 -(1+v(z)^2)d(p',p)
 \leq v(z)(p'(z)-p(z))
 \leq (1+v(z)^2)d(p',p).
\]
Summing over \(z\) places \(V_v(p')-V_v(p)\) between
\(-Md(p',p)\) and \(Md(p',p)\). Suppose now
\(V_v(p)>V_v(r)\), and write
\(\delta=V_v(p)-V_v(r)>0\). Choose
\(\varepsilon=\delta/(4M)\). If the relative lottery distances
from \(p'\) to \(p\) and from \(r'\) to \(r\) are both
less than \(\varepsilon\), then
\[
 V_v(p')-V_v(r')>
 \delta-2M\varepsilon=\delta/2>0.
\]
This proves the actual neighborhood continuity used by the referenced
finite-prize theorem, including boundary lotteries and constant values
(for which there is no strict comparison to check).

For \(w(z)=av(z)+b\), with \(a>0\), summation and normalization give
\[
 V_w(p)=\sum_zp(z)(av(z)+b)
       =aV_v(p)+b\sum_zp(z)=aV_v(p)+b.
\]
Subtracting the same expression for another lottery preserves the sign
of every comparison. This implication also covers a singleton prize
set and all-indifferent values. It does not prove the converse uniqueness
claim for arbitrary domains or permit negative scaling.

For the editorial four-profile value, cooperation produces the sole
value-one profile with probability \(q\) and a value-zero profile
otherwise. Its expected value is therefore \(q\). Defection never
produces \(\langle C,C\rangle\), so its expected value is zero.
The strict and tied comparisons follow respectively when \(q>0\)
and \(q=0\). If the player mixes, cooperating with probability
\(t\), independent draws produce value \(tq\); hence pure
cooperation maximizes when \(q>0\), while every own mixture ties
when \(q=0\). The construction specifies a coherent expected-value
preference; it neither proves that a real participant has it nor
justifies assigning this structural value normatively.
\end{proof}
\end{annexitem}
